\documentclass[11pt]{article}
\usepackage[margin=1in]{geometry}
\usepackage{amsmath,amssymb}
\usepackage[T1]{fontenc}
\usepackage{lmodern}
\usepackage{hyperref}
\usepackage{microtype}

\title{Hybrid Retrieval for Code Chunk Selection in the TBYC Evaluation Pipeline}
\author{}
\date{\today}

\begin{document}
\maketitle

\begin{abstract}
This note documents the retrieval architecture used to select repository code chunks for issue-conditioned reasoning in the TBYC evaluation pipeline. The implemented system is a hybrid sparse--dense retriever that combines lexical matching (BM25), semantic similarity (CodeBERT embeddings indexed with FAISS), and rank-level aggregation (Reciprocal Rank Fusion, RRF). We focus on the high-value modeling decisions and scoring equations that govern retrieval quality.
\end{abstract}

\section{Problem Formulation}
Given an issue query $q$ (constructed from issue title and body) and a corpus of code chunks $\mathcal{D}=\{d_1,\dots,d_N\}$, the goal is to return a ranked subset of chunks most relevant for downstream artifact extraction and thought generation.

The pipeline instantiates two complementary retrieval signals:
\begin{enumerate}
  \item a sparse lexical relevance score $s_{\text{bm25}}(d,q)$,
  \item a dense semantic similarity score $s_{\text{dense}}(d,q)$.
\end{enumerate}
Final ranking is obtained by fusing their rank positions rather than directly mixing raw score magnitudes.

\section{Sparse Retrieval: BM25}
Each chunk is tokenized with a code-aware regex tokenizer and indexed using BM25. The scoring function is the standard Okapi BM25 form:
\begin{equation}
\label{eq:bm25}
s_{\text{bm25}}(d,q)
= \sum_{t \in q}
\operatorname{IDF}(t)
\cdot
\frac{f(t,d)(k_1+1)}{f(t,d)+k_1\left(1-b+b\frac{|d|}{\operatorname{avgdl}}\right)}.
\end{equation}

Here $f(t,d)$ is term frequency in chunk $d$, $|d|$ is chunk length, and $\operatorname{avgdl}$ is average chunk length. BM25 contributes robust lexical recall for exact identifiers, symbol names, and issue-specific terminology.

\section{Dense Retrieval: CodeBERT Embedding Similarity}
The dense branch encodes both queries and chunks with a transformer encoder (CodeBERT family). For an input sequence, the model output is mean-pooled over non-padding token states and then L2-normalized:
\begin{equation}
\hat{\mathbf{e}}(x)
= \frac{\operatorname{Pool}(x)}{\lVert \operatorname{Pool}(x) \rVert_2}.
\end{equation}

Chunk vectors are precomputed and stored; retrieval uses FAISS \texttt{IndexFlatIP} for maximum inner-product search. Because vectors are normalized, inner product equals cosine similarity:
\begin{equation}
\label{eq:dense}
s_{\text{dense}}(d,q)
= \langle \hat{\mathbf{e}}(d), \hat{\mathbf{e}}(q) \rangle
= \cos\big(\hat{\mathbf{e}}(d),\hat{\mathbf{e}}(q)\big).
\end{equation}

This branch improves retrieval when lexical overlap is weak but semantic intent is aligned.

\section{Rank Fusion: Reciprocal Rank Fusion (RRF)}
Rather than calibrating BM25 and dense scores into a common numeric scale, the system fuses ranked lists via RRF. Let $\operatorname{rank}_r(d)$ denote the 1-based position of chunk $d$ in retriever $r \in \{\text{bm25},\text{dense}\}$. The fused score is
\begin{equation}
\label{eq:rrf}
\operatorname{RRF}(d)
= \sum_{r \in \{\text{bm25},\text{dense}\}}
\frac{1}{k + \operatorname{rank}_r(d)}.
\end{equation}

The implementation uses $k=60$ by default. RRF is particularly suitable in hybrid retrieval because it is rank-robust, avoids brittle score normalization, and rewards cross-retriever consensus.

\section{Why This Hybrid Works for Code Retrieval}
The architecture is intentionally complementary:
\begin{itemize}
  \item \textbf{BM25} is strong on literal matches (APIs, symbols, identifiers, error strings).
  \item \textbf{CodeBERT dense search} captures semantic relatedness across paraphrase and abstraction mismatch.
  \item \textbf{RRF} provides a stable, scale-invariant aggregation layer.
\end{itemize}

In practice, this reduces failure modes of single-signal retrieval: sparse-only methods underperform when wording diverges, while dense-only methods may miss precise lexical anchors critical in code search.

\section{Implementation Anchors}
The mathematical pipeline above maps directly to the codebase:
\begin{itemize}
  \item BM25 construction and query scoring in \texttt{src/tbyc\_dataset/evaluation/pipeline.py}.
  \item Dense embedding computation, normalization, and FAISS index search in the same module.
  \item RRF aggregation via the retrieval fusion utility in that module.
\end{itemize}

\section{Conclusion}
The retrieval subsystem is a principled hybrid RAG-style ranker for code chunks: lexical relevance via BM25 (Eq.~\ref{eq:bm25}), semantic similarity via normalized CodeBERT embeddings and FAISS inner product (Eq.~\ref{eq:dense}), and robust rank aggregation via RRF (Eq.~\ref{eq:rrf}). This design prioritizes retrieval fidelity for issue-grounded code understanding without relying on fragile score-space calibration.

\end{document}
